%%%%%%%%%%%%%%%%%%%%%%%%%%%%%%%%%%%%%%%%%%%%%%%%%%%%%%%%%%%%%%%%%%%%%
%  GSoC 2026 Proposal — robstatpy: Python Wrappers for RobStatTM   %
%  Author : Aakarsh Gupta                                   
%%%%%%%%%%%%%%%%%%%%%%%%%%%%%%%%%%%%%%%%%%%%%%%%%%%%%%%%%%%%%%%%%%%%%

\documentclass[11pt, a4paper]{article}

%--------------------------------------------------------------------
% Packages
%--------------------------------------------------------------------
\usepackage[margin=1in]{geometry}
\usepackage{graphicx}
\usepackage{xcolor}
\usepackage{hyperref}
\usepackage{titlesec}
\usepackage{enumitem}
\usepackage{booktabs}
\usepackage{array}
\usepackage{tabularx}
\usepackage{fancyhdr}
\usepackage{amsmath}
\usepackage{amssymb}
\usepackage{listings}
\usepackage{microtype}
\usepackage{parskip}
\usepackage{mdframed}
\usepackage[T1]{fontenc}
\usepackage{lmodern}

%--------------------------------------------------------------------
% Colors
%--------------------------------------------------------------------
\definecolor{uwpurple}{HTML}{4B2E83}
\definecolor{uwgold}{HTML}{85754D}
\definecolor{accentblue}{HTML}{1A73E8}
\definecolor{lightgray}{HTML}{F5F5F5}
\definecolor{darkgray}{HTML}{2E2E2E}
\definecolor{gsocred}{HTML}{DB4437}
\definecolor{softyellow}{HTML}{FFF8E1}
\definecolor{yellowborder}{HTML}{F9A825}

%--------------------------------------------------------------------
% Hyperref
%--------------------------------------------------------------------
\hypersetup{
  colorlinks  = true,
  linkcolor   = uwpurple,
  urlcolor    = accentblue,
  citecolor   = accentblue,
  pdftitle    = {GSoC 2026: robstatpy, Aakarsh Gupta},
  pdfauthor   = {Aakarsh Gupta},
}

%--------------------------------------------------------------------
% Section styling
%--------------------------------------------------------------------
\titleformat{\section}
  {\large\bfseries\color{uwpurple}}
  {\thesection.}{0.6em}{}
  [\vspace{-0.3em}\textcolor{uwgold}{\titlerule[1.2pt]}]

\titleformat{\subsection}
  {\normalsize\bfseries\color{black}}
  {\thesubsection.}{0.5em}{}

\titlespacing{\section}{0pt}{1.6em}{0.7em}
\titlespacing{\subsection}{0pt}{0.9em}{0.3em}

%--------------------------------------------------------------------
% Header / Footer
%--------------------------------------------------------------------
\pagestyle{fancy}
\fancyhf{}
\renewcommand{\headrulewidth}{0.4pt}
\renewcommand{\headrule}{\hbox to\headwidth{%
  \color{uwpurple}\leaders\hrule height \headrulewidth\hfill}}
\fancyhead[L]{\small\color{uwpurple}\textbf{GSoC 2026}
              \textcolor{black}{$\mid$ \texttt{robstatpy}}}
\fancyhead[R]{\small\color{black}Aakarsh Gupta}
\fancyfoot[C]{\small\color{black}\thepage}

%--------------------------------------------------------------------
% Code listing style
%--------------------------------------------------------------------
\lstdefinestyle{python}{
  language        = Python,
  basicstyle      = \ttfamily\footnotesize,
  keywordstyle    = \bfseries\color{accentblue},
  commentstyle    = \itshape\color{darkgray},
  stringstyle     = \color{gsocred},
  numberstyle     = \tiny\color{darkgray},
  numbers         = left,
  numbersep       = 8pt,
  frame           = single,
  framerule       = 0.4pt,
  rulecolor       = \color{uwpurple!40},
  backgroundcolor = \color{lightgray},
  breaklines      = true,
  showstringspaces= false,
  tabsize         = 4,
}
\lstdefinestyle{rstyle}{
  language        = R,
  basicstyle      = \ttfamily\footnotesize,
  keywordstyle    = \bfseries\color{accentblue},
  commentstyle    = \itshape\color{darkgray},
  stringstyle     = \color{gsocred},
  numberstyle     = \tiny\color{darkgray},
  numbers         = left,
  numbersep       = 8pt,
  frame           = single,
  framerule       = 0.4pt,
  rulecolor       = \color{uwpurple!40},
  backgroundcolor = \color{lightgray},
  breaklines      = true,
  showstringspaces= false,
  tabsize         = 4,
}
\lstset{style=python}

%--------------------------------------------------------------------
% Custom environments
%--------------------------------------------------------------------
\newmdenv[
  backgroundcolor  = softyellow,
  linecolor        = yellowborder,
  linewidth        = 1pt,
  innerleftmargin  = 10pt,
  innerrightmargin = 10pt,
  innertopmargin   = 7pt,
  innerbottommargin= 7pt,
  skipabove        = 4pt,
  skipbelow        = 4pt,
]{todobox}

%====================================================================
%  DOCUMENT
%====================================================================
\begin{document}

%--------------------------------------------------------------------
% TITLE BLOCK
%--------------------------------------------------------------------
\thispagestyle{empty}
\begin{center}
  {\LARGE\bfseries\color{uwpurple}
    \texttt{robstatpy}: Python Wrappers for\\[5pt]
    Modern Robust Statistical Estimation
  }\\[10pt]
  {\large\color{black}
    Google Summer of Code 2026, Project Proposal
  }\\[5pt]
  {\normalsize\color{black}
    Organization: \textbf{R Project for Statistical Computing}
    \quad $\cdot$ \quad
    Submitted: 
  }
\end{center}

\vspace{0.5em}
\noindent\textcolor{uwpurple}{\rule{\linewidth}{2pt}}
\vspace{0.3em}

%--------------------------------------------------------------------
% SECTION 1 — PROJECT TITLE & SHORT DESCRIPTION
%--------------------------------------------------------------------
\section{Project Title and Short Description}

\subsection*{Title}

\textbf{\texttt{robstatpy}: Python Wrappers for the \texttt{RobStatTM}
Robust Statistics R~Package}

\subsection*{Short Description}

The \texttt{RobStatTM} R package (Maronna, Martin, Yohai \& Salibian-Barrera, 2019)
implements state-of-the-art robust estimators, including MM-regression,
Rocke's S-estimator for covariance matrices, robust PCA via M-scale
minimization, and Bianco--Yohai logistic regression, none of which have a
direct equivalent in the Python ecosystem.
This project will produce \textbf{\texttt{robstatpy}}, a Python library that
makes the full range of \texttt{RobStatTM} functionality available to Python users
through lightweight \texttt{rpy2}
wrappers, delivering a NumPy/pandas-native interface whose numerical outputs
and graphical analytics reproduce the R originals.
Every wrapper will ship with Python-style documentation modelled on the
\texttt{RobStatTM} R man pages, a \texttt{pytest} validation suite, and
reproducible Jupyter notebook examples.

%--------------------------------------------------------------------
% SECTION 2 — STUDENT INFORMATION
%--------------------------------------------------------------------
\section{Student Information}

\subsection{Contact Details}

\renewcommand{\arraystretch}{1.25}
\begin{tabularx}{\linewidth}{@{} >{\bfseries}l X @{}}
  \toprule
  Field            & Details \\
  \midrule
  Full Name        & Aakarsh Gupta \\
  Mailing Address  & Alta 101, 4735 21st Ave NE, Seattle, WA 98105 \\
  Phone            & +1\,(313)\,663-0680 \\
  University Email & \href{mailto:aakarshg@uw.edu}{\texttt{aakarshg@uw.edu}} \\
  Personal Email   & \href{mailto:aakarshgupta919@gmail.com}
                        {\texttt{aakarshgupta919@gmail.com}} \\
  GitHub           & \href{https://github.com/Aakarsh751}
                        {\texttt{github.com/Aakarsh751}} \\
  Project Repo     & \href{https://github.com/Aakarsh751/robstattm-pyport-AG-prop}
                        {\texttt{github.com/Aakarsh751/robstattm-pyport-AG-prop}} \\
  LinkedIn         & \href{https://linkedin.com/in/aakarsh-gupta-128204218/}
                        {\texttt{linkedin.com/in/aakarsh-gupta-128204218}} \\
  Communication    & Zoom, Google Meet \\
  Time Zone        & Pacific Time (PT,\;UTC$-$7) \\
  \bottomrule
\end{tabularx}
\renewcommand{\arraystretch}{1}

\subsection{Academic Affiliation}

\begin{tabularx}{\linewidth}{@{} l X r @{}}
  \toprule
  \textbf{Institution} & \textbf{Degree} & \textbf{Period} \\
  \midrule
  University of Washington, Seattle
    & B.S.\ in Computational Finance and Risk Management (CFRM)
    & Aug 2025 -- Mar 2027 \\[4pt]
  University of Michigan
    & B.S.\ in Data Science \textit{(Dean's List; International Undergraduate Scholar)}
    & Aug 2023 -- Aug 2025 \\
  \bottomrule
\end{tabularx}

\vspace{0.6em}
\noindent
I am currently a \textbf{junior} with an expected graduation of \textbf{March~2027}.
I began my undergraduate studies in Data Science at the University of Michigan,
where I was placed on the Dean's List and was awarded the International Undergraduate
Scholarship.
In August~2025, I transferred to the University of Washington to pursue a
B.S.\ in Computational Finance and Risk Management, a programme that sits at the
intersection of mathematical statistics, numerical computation, and financial
theory, which is exactly the disciplinary blend required to bridge rigorous
robust-statistics research and the Python data-science ecosystem.

\subsection{Relevant Coursework}

The following UW courses form the academic backbone of this project:

\vspace{0.3em}
\begin{tabularx}{\linewidth}{@{} l X @{}}
  \toprule
  \textbf{Course} & \textbf{Title and Relevance} \\
  \midrule
  \multicolumn{2}{@{}l}{\textit{Completed}} \\[2pt]
  CFRM 425 & \textit{R Programming for Quantitative Finance}:
             hands-on R package development and financial modelling in R \\
  CFRM 410 & \textit{Probability \& Statistics for Computational Finance}:
             estimation theory, asymptotic methods, sampling distributions \\
  CFRM 405 & \textit{Mathematical Methods for Quantitative Finance}:
             optimisation, convexity, and numerical root-finding \\
  AMATH 482 & \textit{Computational Methods for Data Analysis}:
              dimensionality reduction, PCA, and spectral methods in Python \\
  AMATH 352 & \textit{Applied Linear Algebra \& Numerical Analysis}:
              matrix decompositions and numerical stability \\[4pt]
  \multicolumn{2}{@{}l}{\textit{Spring 2026 (in progress)}} \\[2pt]
  CFRM 421 & \textit{Machine Learning with Application to Finance}:
             Python-based ML methods applied to financial data, directly
             mirroring the use-case of \texttt{robstatpy} \\
  CFRM 586 & \textit{Financial Time Series Forecasting Methods}:
             modelling heavy-tailed financial returns, a primary application
             domain for robust estimation \\
  \bottomrule
\end{tabularx}

\subsection{Technical Skills}

\begin{itemize}[leftmargin=1.5em, itemsep=3pt]
  \item \textbf{Programming languages:} Python, R, SQL, C++
  \item \textbf{Python ecosystem:}
        NumPy, pandas, scikit-learn, matplotlib, PyTorch, TensorFlow,
        Selenium; REST API development with FastAPI and Streamlit
  \item \textbf{R ecosystem:}
        \texttt{ggplot2}, \texttt{tidyr}, \texttt{dplyr},
        \texttt{RobStatTM}, \texttt{robustbase};
        DataCamp-certified in R~Programming, Data~Manipulation with dplyr,
        and Data~Visualisation with ggplot2
  \item \textbf{Development tools:}
        Git/GitHub, Jupyter, RStudio, Docker, Anaconda
  \item \textbf{Finance:}
        Bloomberg Finance Fundamentals; quantitative equity and
        crypto-market analysis; portfolio risk modelling
\end{itemize}

\subsection{Relevant Experience and Projects}

\begin{description}[leftmargin=0pt, itemsep=8pt]

  \item[\normalfont\textbf{IT Risk Management and Analyst Intern,
        Virginia Department of Motor Vehicles}
        ]
  Spearheaded a proof-of-concept initiative to evaluate AI and NLP for
  analysing call-centre interactions, assessing customer sentiment,
  call-resolution efficiency, agent performance, and fraud-related risk
  indicators.
  Built FastAPI back-ends and Streamlit front-ends to deploy and interact
  with large language models, enabling real-time inference and operational
  decision support.
  Designed and automated intelligent workflows using n8n to streamline
  repetitive tasks, and processed unstructured speech data to produce
  cost-benefit analyses guiding in-house versus outsourced AI deployment
  decisions.
  This role directly exercised the API-design, packaging, and data-pipeline
  skills that are central to \texttt{robstatpy}.

  \item[\normalfont\textbf{Undergraduate Research Volunteer,
        University of Michigan -- Imitation Learning}
        ]
  Conducted research on Imitation Learning under faculty guidance, training
  reinforcement-learning and human-guided agents with OpenAI Gym.
  Implemented the DAGGER and HG-DAGGER algorithms from the primary literature
  to improve agent performance, adaptability, and learning efficiency, and
  evaluated all models rigorously using PyTorch.
  This experience demonstrates the ability to read formal algorithm
  specifications, translate them into well-tested, reproducible code,
  and verify results against published benchmarks, a skill directly
  applicable to wrapping the MM- and S-estimation routines in
  \texttt{RobStatTM}.

  \item[\normalfont\textbf{Deep Learning Algorithm for Asthenopia Detection}
        ]
  Independently conceived and developed a deep-learning image-classification
  pipeline for the early detection of asthenopia (visual fatigue).
  The project was recognised with a \textbf{Gold Award at the Indian Science
  and Engineering Fair} and achieved a top-100 national selection, providing
  evidence of the capacity to drive an independent research project from
  initial conception through to a peer-recognised outcome.

\end{description}

\subsection{Why My Background Aligns with This Project}

My training in Computational Finance and Risk Management (CFRM) is
uniquely suited to this project.
Robust statistics is not an abstract topic for me: it is central to the
quantitative methods I am now beginning to study and apply.
In CFRM~410 (Probability and Statistics), I learned estimation
theory and asymptotic analysis that will allow me to understand the underpinnings of robust MM-estimators and
S-Estimators that are covered in Maronna et al.\ (2019).
In CFRM~425 (R Programming), I developed R packages and gained hands-on
familiarity with the object system that \texttt{RobStatTM} uses
internally.
In CFRM~405 (Mathematical Methods), I studied the optimisation
numerical algorithms that drive iterative robust estimators.
AMATH~482 (Computational Methods for Data Analysis) gave me direct
experience implementing PCA and dimensionality reduction in Python, the
exact techniques that \texttt{pcaRobS} robustifies.

This combination means I understand \emph{both}
the R code base and its statistical foundations, and the Python
ecosystem that \texttt{robstatpy} targets.
I am not just a software engineer wrapping an unfamiliar library; I have
the theoretical knowledge to verify correctness, interpret results, and
extend the project into downstream applications such as robust portfolio
construction and risk analysis (RPCRA).

\subsection{Relationship with Mentors}

I have been in active communication with \textbf{Prof.\ R.\ Douglas Martin}
since \textbf{January 2026}, attending weekly meetings to discuss the project
scope, the Python API design, and the technical challenges of bridging R's S3
object system with Python's language model.
These conversations have directly shaped the implementation strategy presented
in this proposal.
I look forward to further engaging with the full mentorship team,
Prof.\ Matias Salibian-Barrera and Brian G.\ Peterson, as the project progresses.

\subsection{Mentors}

\renewcommand{\arraystretch}{1.2}
\begin{tabularx}{\linewidth}{@{} l l X @{}}
  \toprule
  \textbf{Name} & \textbf{Role} & \textbf{Contact} \\
  \midrule
  Prof.\ R.\ Douglas Martin
    & Primary mentor
    & \href{mailto:martinrd3d@gmail.com}{\texttt{martinrd3d@gmail.com}} \\
  Prof.\ Matias Salibian-Barrera
    & Co-mentor
    & \href{mailto:matias@stat.ubc.ca}{\texttt{matias@stat.ubc.ca}} \\
  Brian G.\ Peterson
    & Co-mentor
    & \href{mailto:brian@braverock.com}{\texttt{brian@braverock.com}} \\
  \bottomrule
\end{tabularx}
\renewcommand{\arraystretch}{1}

\subsection{Time Commitment and Availability}

I am applying for the \textbf{large project size (350~hours)}.
I am fully dedicated and available for the entire duration of the GSoC
coding period (late May through mid-November 2026), with
\textbf{30--35 hours per week} committed to \texttt{robstatpy}
development.
I will continue attending weekly meetings with Prof.\ Martin and will be
responsive to all three mentors via email, Zoom, and GitHub.

%--------------------------------------------------------------------
% SECTION 3 — BENEFITS TO COMMUNITY
%--------------------------------------------------------------------
\section{Benefits to the Community}

The robust-statistics methods described in \texttt{RobStatTM}
including MM-regression, Rocke's S-estimator, robust PCA via M-scale minimization, among others,
are the algorithmic companions to \emph{Robust Statistics: Theory and
Methods} (Maronna, Martin, Yohai \& Salibian-Barrera, 2nd~ed., 2019),
a leading modern textbook in the field.
Today, a Python practitioner who wants to reproduce a textbook example or
apply these methods in their research or applications faces two unappealing choices:
install and learn R, or re-implement algorithms from scratch in Python.

\texttt{robstatpy} removes this barrier by providing a
pip\footnote{\texttt{pip} is the standard Python package installer.
Running \texttt{pip install robstatpy} downloads and installs the
package and all its dependencies automatically.}-installable
Python package that:

\begin{enumerate}[leftmargin=2em]
  \item \textbf{Exposes the full \texttt{RobStatTM} API} through
    \texttt{rpy2} wrappers, accepting NumPy arrays and pandas DataFrames
    and returning native Python objects.
  \item \textbf{Guarantees numerical fidelity:} every output agrees
    with the R original to machine precision, because the R code itself
    executes underneath.
  \item \textbf{Lowers the maintenance burden:} because the wrappers
    delegate to R, any upstream improvement to \texttt{RobStatTM} is
    automatically available in Python without duplicating development
    effort.
  \item \textbf{Reproduces figures and numerical summaries} using
    \texttt{matplotlib} and \texttt{plotnine}\footnote{\texttt{plotnine}
    is a Python implementation of R's \texttt{ggplot2}
    grammar-of-graphics, enabling near-identical plot syntax in Python.
    See \url{https://plotnine.org/} and \url{https://realpython.com/ggplot-python/}.}
    so that Python users can produce visualisations that closely
    mirror the R originals.
  \item \textbf{Provides a comparative study} of wrapper-based versus
    direct-porting strategies, producing benchmarks and guidelines that
    will inform future R-to-Python interoperability projects.
\end{enumerate}

\noindent
\textbf{Who benefits:}
\begin{itemize}[leftmargin=1.5em, itemsep=2pt]
  \item \emph{Data scientists, ML engineers, and Python-only practitioners}
    who need a rich toolkit of modern robust estimation methods without
    learning R. \texttt{robstatpy} provides a comprehensive, self-contained
    Python library: users call familiar Python functions with NumPy arrays
    and pandas DataFrames and never need to write, read, or understand any
    R code.
  \item \emph{Students and researchers} who study the Maronna~et~al.\
    textbook and want to apply the methods in Jupyter rather than RStudio.
\end{itemize}

%--------------------------------------------------------------------
% SECTION 4 — RELATED WORK
% SECTION 4 — RELATED WORK
%--------------------------------------------------------------------
\section{Related Work}

\subsection{Existing Python Robust-Statistics Packages}

\begin{itemize}[leftmargin=1.5em, itemsep=4pt]
  \item \textbf{scikit-learn} provides \texttt{MinCovDet} (fast MCD) and
    \texttt{HuberRegressor}, but lacks MM-regression, S-estimators for
    covariance, and robust PCA beyond standard PCA.
  \item \textbf{statsmodels} offers \texttt{RLM} (M-regression with
    Huber or Tukey bisquare weights) but does not implement the
    two-stage S+M refinement that defines MM-estimation, nor
    \texttt{lmrobdet.control}-style tuning.
  \item \textbf{RobPy} (Leyder et al., 2024; arXiv:2411.01954) is a
    comprehensive Python robust-statistics package offering MM-regression,
    S-regression, fast LTS, MCD-based covariance estimators (including
    cellwise MCD), ROBPCA, and cellwise outlier detection, all built as
    native re-implementations drawing from \texttt{robustbase},
    \texttt{rrcov}, and \texttt{cellWise}. However, it does not expose
    the \texttt{RobStatTM} function set: the DCML regression family,
    robust logistic regression, additional multivariate scatter estimators,
    and the SM-PCA variant, along with the full
    \texttt{lmrobdetDCML}/\texttt{lmrobdetMM} tuning infrastructure
    codified in Maronna et al.\ (2019), remain without a Python equivalent.
  \item \textbf{robustats} (PyPI) provides C-backed univariate
    estimators (median, MAD, Qn) but nothing multivariate.
\end{itemize}

%--------------------------------------------------------------------
% SECTION 5 -- PROPOSED WORK
%--------------------------------------------------------------------
\section{Proposed Work}

\texttt{robstatpy} pursues a clear, phased strategy: first expose the
complete \texttt{RobStatTM} function set to Python users via lightweight
\texttt{rpy2} wrappers for immediate, numerically exact access; then
provide alternative approaches for flexibility and portability.
This section describes the primary wrapper approach, presents code
examples for key functions already implemented, and reports validation
results confirming numerical equivalence with R.

\subsection{Why Wrappers Are the Right First Step}

The \texttt{RobStatTM} package is \textbf{actively maintained} by its
original authors, the same researchers who proved the theoretical
properties of these estimators.
A wrapper-based approach has three decisive advantages over full
re-implementation:

\begin{enumerate}[leftmargin=2em]
  \item \textbf{Exact numerical equivalence.}  Because the R code
    executes underneath, Python users get bit-identical results.
    This is essential for textbook reproducibility and for
    cross-validating future pure-Python ports.
  \item \textbf{Zero duplication of maintenance.}  When the R authors
    fix a bug or add an estimator, the Python wrapper inherits the
    improvement immediately.
    A pure-Python re-implementation creates a second code base that
    must be independently maintained and validated.
  \item \textbf{Speed of delivery.}  Wrapping a function requires
    writing conversion glue, tests, and documentation, not
    re-deriving the algorithm.
    This makes it feasible to expose the \emph{entire} API within a
    single GSoC period, rather than a subset.
\end{enumerate}

\noindent
The project will \emph{also} produce pure-Python re-implementations,
beginning with \texttt{locScaleM} and \texttt{scaleM} and extending
toward a complete native Python port of the \texttt{RobStatTM} library.
The wrapper layer serves as the validated ground truth against which
every pure-Python function is tested: once a ported function passes
numerical equivalence checks against the rpy2 wrapper, it can replace
the wrapper call, progressively removing the R dependency.
This dual-track strategy, wrappers first for correctness and coverage,
then native ports for independence and performance, provides the most
robust path to a fully self-contained Python library while ensuring that
no accuracy is sacrificed along the way.

\subsection{Strategy~A: \texttt{rpy2} Wrapper Layer (Primary Approach)}

Every public \texttt{RobStatTM} function will be wrapped in a Python function
that:
\begin{enumerate}[leftmargin=2em]
  \item Accepts NumPy arrays and/or pandas DataFrames.
  \item Converts inputs to R objects via \texttt{rpy2.robjects} and
    \texttt{numpy2ri}/\texttt{pandas2ri} converters.
  \item Calls the R function inside \texttt{RobStatTM}.
  \item Extracts named elements from the returned R list using
    \texttt{.rx2()}.
  \item Returns a Python dataclass (or named dict) whose fields match
    the R return values.
\end{enumerate}

\noindent
\textbf{Architecture diagram:}

\begin{center}
\begin{tabular}{ccccc}
  \fbox{\texttt{Python user code}} &
  $\xrightarrow{\text{NumPy/pandas}}$ &
  \fbox{\texttt{robstatpy API}} &
  $\xrightarrow{\text{rpy2 bridge}}$ &
  \fbox{\texttt{RobStatTM (R)}} \\[6pt]
  & & $\downarrow$ & & $\downarrow$ \\[4pt]
  & & \fbox{\texttt{Python dataclass}} &
  $\xleftarrow{\text{.rx2() extract}}$ &
  \fbox{\texttt{R list}} \\
\end{tabular}
\end{center}

\noindent\textbf{Key design decisions:}

\begin{itemize}[leftmargin=1.5em, itemsep=3pt]
  \item \textbf{Global conversion context.}
    \texttt{rpy2}~3.6+ stores conversion rules in a
    \texttt{contextvars.ContextVar}, which is lost across Jupyter's
    async task boundaries.
    We patch this by calling \texttt{set\_conversion(default\_converter)}
    at import time, ensuring conversions persist across cells.
  \item \textbf{R dot-to-Python underscore mapping.}
    \texttt{importr()} automatically maps \texttt{lmrobdet.control} to
    \texttt{lmrobdet\_control}.
    We document both names in every docstring.
  \item \textbf{Lazy R package loading.}
    \texttt{RobStatTM} is loaded via \texttt{importr} only on first
    use, so importing \texttt{robstatpy} in a pure-Python context
    (e.g.\ for type hints or documentation builds) does not require R.
\end{itemize}


\noindent
The complete \texttt{locScaleM} and \texttt{scaleM} wrapper implementations,
pure-Python re-implementations, and detailed numerical test results are
provided in the companion notebook in the project repository:
\href{https://github.com/Aakarsh751/robstattm-pyport-AG-prop/blob/main/robstattm/python/robstatpy\_comparison\_rpy2.ipynb}%
{\texttt{robstatpy\_comparison\_rpy2.ipynb}}.

\subsubsection*{Code Example: \texttt{rpy2} Wrapper for \texttt{lmrobdetMM}}

\begin{lstlisting}[style=python, caption={rpy2 wrapper for lmrobdetMM (regression; demonstrated with mopt family, efficiency=0.95)}]
def lmrobdet_mm(formula, data, family="mopt", efficiency=0.95):
    """Robust MM-regression via RobStatTM::lmrobdetMM().

    Parameters
    ----------
    formula : str
        R-style formula, e.g. 'zinc ~ copper'.
    data : pd.DataFrame
        Data frame (converted to R via pandas2ri).
    family : str
        Psi function family for IRLS refinement.
    efficiency : float
        Target asymptotic efficiency.

    Returns
    -------
    dict with 'coefficients', 'scale', 'residuals',
         'fitted_values', 'rweights', 'r_squared'.
    """
    from rpy2.robjects import pandas2ri
    pandas2ri.activate()
    r_df = ro.conversion.py2rpy(data)
    ctrl = robstattm.lmrobdet_control(
        family=family, efficiency=efficiency)
    fit = robstattm.lmrobdetMM(
        ro.Formula(formula), data=r_df, control=ctrl)
    return {
        "coefficients":  np.array(fit.rx2("coefficients")),
        "scale":         fit.rx2("scale")[0],
        "residuals":     np.array(fit.rx2("residuals")),
        "fitted_values": np.array(fit.rx2("fitted.values")),
        "rweights":      np.array(fit.rx2("rweights")),
        "r_squared":     fit.rx2("r.squared")[0],
    }
\end{lstlisting}


\subsection{Validation Results}

The wrapper code examples above (\texttt{lmrobdet\_mm} and
\texttt{cov\_rob\_mm}) are representative of the full wrapper layer.
Each wrapper calls the corresponding \texttt{RobStatTM} R function
via \texttt{rpy2} and returns Python-native objects; because R
executes underneath, outputs are numerically identical to a direct R
session.
Full wrapper code and numerical result tables for all five functions
(\texttt{locScaleM}, \texttt{scaleM}, \texttt{lmrobdetMM},
\texttt{covRobMM}, \texttt{covRobRocke}, \texttt{pcaRobS}) are
provided in the companion test document:
\href{https://github.com/Aakarsh751/robstattm-pyport-AG-prop/blob/main/docs/gsoc2026_proposal/robstatpy_tests.pdf}%
{\texttt{robstatpy\_tests.pdf}}.
The rpy2 wrapper notebook (14/14 checks pass) is at
\href{https://github.com/Aakarsh751/robstattm-pyport-AG-prop/blob/main/robstattm/python/robstatpy\_comparison\_rpy2.ipynb}%
{\texttt{robstatpy\_comparison\_rpy2.ipynb}}.

\subsection{Additional Wrapped Functions (via rpy2)}

The rpy2 notebook validates all five core \texttt{RobStatTM} functions
by calling R and checking outputs against expected values:

\renewcommand{\arraystretch}{1.2}
\begin{tabularx}{\linewidth}{@{} l l X c @{}}
  \toprule
  \textbf{Function} & \textbf{Dataset} & \textbf{Validation} & \textbf{Status} \\
  \midrule
  \texttt{lmrobdetMM}  & mineral & Coefficients, scale, residuals, rweights & PASS \\
  \texttt{covRobMM}    & wine    & Center vector, covariance matrix, distances & PASS \\
  \texttt{covRobRocke} & wine    & Center vector, distances & PASS \\
  \texttt{pcaRobS}     & bus     & Proportion of variance, eigenvectors & PASS \\
  \bottomrule
\end{tabularx}
\renewcommand{\arraystretch}{1}

\vspace{0.5em}
\noindent\textbf{Selected numerical outputs --- pure R vs.\ \texttt{robstatpy} (rpy2 wrapper):}

\renewcommand{\arraystretch}{1.2}
\begin{tabularx}{\linewidth}{@{} l l r r @{}}
  \toprule
  \textbf{Function} & \textbf{Output Field}
    & \textbf{Pure R} & \textbf{robstatpy} \\
  \midrule
  \multicolumn{4}{@{}l}{\textit{lmrobdetMM (mineral, mopt, eff\,=\,0.95)}} \\
  & Intercept          & $15.2012$ & $15.2012$ \\
  & Slope (copper)     & $0.0126$  & $0.0126$  \\
  & Scale              & $9.9966$  & $9.9966$  \\
  & Robust $R^2$       & $0.0082$  & $0.0082$  \\[3pt]
  \multicolumn{4}{@{}l}{\textit{covRobMM (wine, 13 variables, \texttt{set.seed(42)})}} \\
  & Center{[}1..5{]} & $13.75,\;1.77,\;2.46,\;17.06,\;106.19$
                     & $13.75,\;1.77,\;2.46,\;17.06,\;106.19$ \\
  & Max Mahalanobis dist & $164.36$ & $164.36$ \\[3pt]
  \multicolumn{4}{@{}l}{\textit{covRobRocke (wine, 13 variables, \texttt{set.seed(42)})}} \\
  & Center{[}1..5{]} & $13.73,\;1.78,\;2.47,\;17.15,\;106.77$
                     & $13.73,\;1.78,\;2.47,\;17.15,\;106.77$ \\
  & Max Mahalanobis dist & $154.15$ & $154.15$ \\[3pt]
  \multicolumn{4}{@{}l}{\textit{pcaRobS (bus, 3 components, \texttt{set.seed(42)})}} \\
  & propSPC{[}1{]} & $0.8611$ & $0.8611$ \\
  & propex         & $0.9731$ & $0.9731$ \\
  \bottomrule
\end{tabularx}
\renewcommand{\arraystretch}{1}

%--------------------------------------------------------------------
% SECTION 6 -- ALTERNATIVE IMPLEMENTATION STRATEGIES
%--------------------------------------------------------------------
\section{Alternative Implementation Strategies}

Two additional approaches were implemented as cross-checks and to
demonstrate different trade-offs in R--Python interoperability.
Both are described briefly here; full code is available in the
project notebooks.

\subsection{Strategy~B: Subprocess + JSON Wrappers}

As a comparison baseline, the project also implements wrappers using
\texttt{subprocess} calls to \texttt{Rscript} with
\texttt{jsonlite}-serialised results.
This approach:
\begin{itemize}[leftmargin=1.5em]
  \item Requires no compiled \texttt{rpy2} dependency (useful on systems
    where building \texttt{rpy2} is difficult).
  \item Incurs higher per-call overhead due to R process startup.
  \item Serves as a cross-check: if both wrappers produce identical
    outputs, we have high confidence in the conversion layer.
\end{itemize}

% [code available in robstatpy\_comparison.ipynb]


\subsection{Strategy~C: Direct Python Re-implementation}

For \texttt{locScaleM} and \texttt{scaleM}, the project provides
pure-Python re-implementations to serve as the GSoC test tasks and to
establish a methodology for future direct porting.

% SECTION 7 -- EVALUATION PLAN
%--------------------------------------------------------------------
\section{Evaluation Plan: Comparative Benchmarks}

A central deliverable is a rigorous comparison of the three implementation
strategies across three dimensions:

\subsection{Metrics}

\begin{enumerate}[leftmargin=2em]
  \item \textbf{Accuracy.}
    For each function, compute $\max_i |y_i^{\text{Py}} - y_i^{\text{R}}|$
    across all return fields, for both the rpy2 wrapper (expected: $0$)
    and the pure-Python port (expected: $< 10^{-6}$).
  \item \textbf{Speed.}
    Wall-clock time per call, measured with \texttt{timeit} over 100
    repetitions on synthetic data of sizes $n \in \{100,\,1000,\,10000\}$
    and $p \in \{5,\,20,\,50\}$.
    Breakdown: R computation time vs.\ conversion overhead.
  \item \textbf{Ergonomics and integration cost.}
    Lines of glue code per function, dependency footprint, and ease of
    installation (does the approach require R to be installed?).
\end{enumerate}

\subsection{Benchmark Protocol}

\begin{itemize}[leftmargin=1.5em, itemsep=3pt]
  \item \textbf{Synthetic data generation:}
    Gaussian base with $\varepsilon = 5\%$ contamination at 10$\sigma$
    to exercise breakdown behaviour.
  \item \textbf{Functions benchmarked:}
    \texttt{locScaleM}, \texttt{scaleM}, \texttt{lmrobdetMM},
    \texttt{covRobMM}, \texttt{covRobRocke}, \texttt{pcaRobS}.
  \item \textbf{Comparison grid:}
    rpy2 wrapper vs.\ subprocess wrapper vs.\ pure-Python port (where
    available) vs.\ nearest \texttt{scikit-learn}/\texttt{statsmodels}
    equivalent (where one exists).
  \item \textbf{Reporting:}
    Results will be presented in a Jupyter notebook with timing tables,
    accuracy heat maps, and commentary.
\end{itemize}

%--------------------------------------------------------------------
% SECTION 8 -- IMPLEMENTATION PLAN
%--------------------------------------------------------------------
\section{Implementation Plan and Timeline}

The project follows a \textbf{22-week large-project timeline}
(350~hours), structured around the GSoC 2026 calendar.

\subsection{Community Bonding Period (May 1 -- May 24)}

\begin{itemize}[leftmargin=1.5em, itemsep=2pt]
  \item Finalise package scaffolding: \texttt{pyproject.toml},
    \texttt{pytest} configuration, CI/CD with GitHub Actions
    (R + Python matrix).
  \item Set up automated testing: R installation in CI, \texttt{rpy2}
    build, dataset fixtures.
  \item Draft the \texttt{robstatpy} API surface with type stubs
    and docstring templates.
  \item Agree on code review cadence and milestone criteria with mentors.
\end{itemize}

\subsection{Phase 1: Core Univariate Wrappers (Weeks 1--3, May 25 -- Jun 14)}

\begin{tabularx}{\linewidth}{@{} l X @{}}
  \toprule
  \textbf{Week} & \textbf{Deliverables} \\
  \midrule
  1 & \texttt{locScaleM} wrapper + pure-Python port + pytest suite
    (bisquare, huber, mopt families; eff $\in \{0.85, 0.90, 0.95\}$).
    Numerical agreement verified to $< 10^{-8}$. \\
  2 & \texttt{scaleM} / \texttt{mscale} wrapper + pure-Python port.
    Subprocess wrapper as cross-check.
    Document relationship between \texttt{scaleM} and \texttt{locScaleM}. \\
  3 & \texttt{lmrobdet.control} wrapper exposing all tuning parameters.
    Comparison notebook: rpy2 vs.\ subprocess vs.\ pure-Python for
    univariate functions. \\
  \bottomrule
\end{tabularx}

\subsection{Phase 2: Robust Regression (Weeks 4--7, Jun 15 -- Jul 12)}

\begin{tabularx}{\linewidth}{@{} l X @{}}
  \toprule
  \textbf{Week} & \textbf{Deliverables} \\
  \midrule
  4 & \texttt{lmrobdetMM} wrapper: formula interface, coefficient
    extraction, residuals, robustness weights, scale. \\
  5 & \texttt{lmrobdetMM} diagnostics: robust $R^2$, fitted values,
    summary table mirroring R's \texttt{summary()} output.
    \texttt{lmrobdetMM.RFPE} wrapper. \\
  6 & \texttt{step.lmrobdetMM} wrapper for stepwise robust regression.
    Reproduce mineral dataset examples (Figs~5.1--5.7). \\
  7 & \textbf{Midterm milestone:} complete regression module with
    full pytest coverage.
    Prepare midterm evaluation report. \\
  \bottomrule
\end{tabularx}

\vspace{0.3em}
\noindent\textbf{Midterm Evaluation (Jul 6--10):}
Deliverables: univariate wrappers, regression wrappers, comparison
notebook, $\geq 90\%$ test coverage on delivered modules.

\subsection{Phase 3: Robust Covariance and PCA (Weeks 8--11, Jul 13 -- Aug 9)}

\begin{tabularx}{\linewidth}{@{} l X @{}}
  \toprule
  \textbf{Week} & \textbf{Deliverables} \\
  \midrule
  8  & \texttt{covRobMM} wrapper: center, covariance, correlation,
     Mahalanobis distances, weights. \\
  9  & \texttt{covRobRocke} wrapper.
     Reproduce wine dataset examples (Fig~6.3). \\
  10 & \texttt{pcaRobS} wrapper: eigenvectors, proportion of variance,
     scores.
     Reproduce bus dataset examples (Fig~6.10). \\
  11 & \texttt{prcompRob} wrapper (convenience wrapper over \texttt{pcaRobS}).
     Distance--distance plot utility function.
     Scree plot utility function. \\
  \bottomrule
\end{tabularx}

\subsection{Phase 4: Benchmarks and Comparative Study (Weeks 12--14, Aug 10 -- Aug 30)}

\begin{tabularx}{\linewidth}{@{} l X @{}}
  \toprule
  \textbf{Week} & \textbf{Deliverables} \\
  \midrule
  12 & Speed benchmarks: rpy2 vs.\ subprocess vs.\ pure-Python vs.\
     scikit-learn/statsmodels baselines.
     Synthetic data generator for contaminated Gaussian. \\
  13 & Accuracy benchmarks: max absolute error across all functions,
     all parameter combinations.
     Tabulate conversion overhead breakdown. \\
  14 & Write-up: comparative analysis Jupyter notebook with tables,
     plots, and recommendations for when each approach is appropriate. \\
  \bottomrule
\end{tabularx}

\subsection{Phase 5: Documentation, Packaging, and Polish (Weeks 15--17, Aug 31 -- Sep 20)}

\begin{tabularx}{\linewidth}{@{} l X @{}}
  \toprule
  \textbf{Week} & \textbf{Deliverables} \\
  \midrule
  15 & Complete API documentation (Sphinx + ReadTheDocs).
     Every function has a docstring modelled on the R man page. \\
  16 & Publish \texttt{robstatpy} on PyPI (sdist + wheel).
     Installation guide covering R dependency setup on
     Linux/macOS/Windows. \\
  17 & Tutorial Jupyter notebooks reproducing key textbook chapters
     (Chapters 2, 5, 6 of Maronna et al.). \\
  \bottomrule
\end{tabularx}

\subsection{Phase 6: PCRA Extension and Final Deliverables (Weeks 18--22, Sep 21 -- Nov 2)}

\begin{tabularx}{\linewidth}{@{} l X @{}}
  \toprule
  \textbf{Week} & \textbf{Deliverables} \\
  \midrule
  18--19 & \textbf{PCRA exploration:}
     Map PCRA's dependency on \texttt{RobStatTM} (calls to
     \texttt{locScaleM}, \texttt{lmrobdetMM} in portfolio risk
     contexts).
     Wrap or port PCRA's core math functions (efficient frontier,
     global minimum variance). \\
  20 & PCRA comparison notebook: R vs.\ Python outputs for Chapter~2
     demo (2-asset frontier, drawdown analysis, risk measures). \\
  21 & Integration testing: end-to-end portfolio analytics workflow
     using \texttt{robstatpy} for robust covariance input. \\
  22 & \textbf{Final submission:}
     Final code review, README polish, contributor guide,
     final evaluation report. \\
  \bottomrule
\end{tabularx}

%--------------------------------------------------------------------
% SECTION 9 -- DELIVERABLES
%--------------------------------------------------------------------
\section{Deliverables}

\subsection{Required Deliverables}

\begin{enumerate}[leftmargin=2em, itemsep=4pt]
  \item \textbf{\texttt{robstatpy} Python package} (pip-installable)
    wrapping the following \texttt{RobStatTM} functions via
    \texttt{rpy2}:
    \begin{itemize}[leftmargin=1.5em]
      \item \texttt{locScaleM}, \texttt{scaleM} (univariate)
      \item \texttt{lmrobdetMM}, \texttt{lmrobdet.control},
        \texttt{lmrobdetMM.RFPE}, \texttt{step.lmrobdetMM}
        (regression)
      \item \texttt{covRobMM}, \texttt{covRobRocke} (covariance)
      \item \texttt{pcaRobS}, \texttt{prcompRob} (PCA)
    \end{itemize}

  \item \textbf{Pure-Python re-implementations} of \texttt{locScaleM}
    and \texttt{scaleM} with verified numerical equivalence, establishing
    the methodology for progressively porting the full library to native
    Python.

  \item \textbf{Pytest test suite} achieving $\geq 90\%$ coverage on
    all wrapper functions, comparing Python outputs against R ground
    truth.

  \item \textbf{Comparative benchmark notebook} evaluating speed,
    accuracy, and ergonomics across rpy2 wrappers, subprocess wrappers,
    and pure-Python ports.

  \item \textbf{Sphinx documentation} published on ReadTheDocs, with
    docstrings modelled on the R man pages.

  \item \textbf{Tutorial Jupyter notebooks} reproducing textbook
    examples from Chapters~2, 5, and~6 of Maronna et al.\ (2019).
\end{enumerate}

\subsection{Stretch Deliverables}

\begin{enumerate}[leftmargin=2em, itemsep=4pt]
  \item \textbf{PCRA Python module} wrapping or porting core portfolio
    analytics functions (\texttt{mathEfrontRiskyMuCov},
    \texttt{mathGmv}, drawdown analysis, risk measures).
  \item \textbf{Additional \texttt{RobStatTM} wrappers:}
    \texttt{lmrobM}, \texttt{BYlogreg} (Bianco--Yohai logistic
    regression), \texttt{rob.linear.test}.
  \item \textbf{Complete native Python port} of the remaining
    \texttt{RobStatTM} functions (\texttt{lmrobdetMM},
    \texttt{covRobMM}, \texttt{covRobRocke}, \texttt{pcaRobS}),
    validated against the rpy2 wrappers, with optional \texttt{numba}
    JIT acceleration for performance-critical inner loops.
\end{enumerate}

%--------------------------------------------------------------------
% SECTION 10 -- DISTRIBUTION
%--------------------------------------------------------------------
\section{Distribution and Packaging}

\texttt{robstatpy} will be distributed through two complementary
channels to maximise accessibility:

\subsection{GitHub Installation (Development and Bleeding-Edge)}

The package will be hosted on GitHub under an open-source license
(MIT or Apache~2.0, to be confirmed with mentors).
Users can install directly from the repository:

\begin{lstlisting}[style=python, caption={GitHub installation}, numbers=none]
pip install git+https://github.com/Aakarsh751/robstatpy.git
\end{lstlisting}

\noindent
This channel provides access to the latest development version and
makes it straightforward for contributors to fork, modify, and submit
pull requests.
The repository will include:
\begin{itemize}[leftmargin=1.5em, itemsep=2pt]
  \item A \texttt{pyproject.toml} with full metadata and dependency
    specifications (\texttt{rpy2 >= 3.5}, \texttt{numpy},
    \texttt{pandas}, \texttt{scipy}, \texttt{matplotlib}).
  \item GitHub Actions CI running the full \texttt{pytest} suite on
    Linux, macOS, and Windows against multiple R versions (4.3, 4.4,
    4.5).
  \item A \texttt{CONTRIBUTING.md} with guidelines for adding new
    wrappers or submitting pure-Python ports.
  \item Tagged releases corresponding to each PyPI publication.
\end{itemize}

\subsection{PyPI Publication (Stable Releases)}

Stable, versioned releases will be published on the Python Package Index
(PyPI), enabling standard installation:

\begin{lstlisting}[style=python, caption={PyPI installation}, numbers=none]
pip install robstatpy
\end{lstlisting}

\noindent
The packaging workflow uses:
\begin{itemize}[leftmargin=1.5em, itemsep=2pt]
  \item \texttt{build} to produce source distributions (\texttt{sdist})
    and pure-Python wheels.
  \item \texttt{twine} for secure upload to PyPI.
  \item Semantic versioning (e.g., \texttt{0.1.0} for the initial GSoC
    release, \texttt{0.2.0} after PCRA extension).
  \item A GitHub Actions release workflow that automatically publishes
    to PyPI when a new git tag is pushed.
\end{itemize}

\noindent
\textbf{R dependency handling:}
Since \texttt{robstatpy} requires a working R installation with
\texttt{RobStatTM}, the package will:
(1) detect \texttt{R\_HOME} automatically at import time,
(2) provide a clear \texttt{ImportError} message with installation
instructions if R is not found, and
(3) include a \texttt{robstatpy.check\_setup()} utility that validates
the R environment and prints version information.
The installation guide will cover R setup on all three major platforms
(Linux via \texttt{apt}/\texttt{conda}, macOS via Homebrew, Windows via
the CRAN installer).

%--------------------------------------------------------------------
% SECTION 11 -- RISKS
%--------------------------------------------------------------------
\section{Risks and Mitigation Strategies}

\renewcommand{\arraystretch}{1.3}
\begin{tabularx}{\linewidth}{@{} >{\raggedright\arraybackslash}p{3.5cm}
                                  >{\raggedright\arraybackslash}X
                                  >{\raggedright\arraybackslash}X @{}}
  \toprule
  \textbf{Risk} & \textbf{Impact} & \textbf{Mitigation} \\
  \midrule
  \texttt{rpy2} build failures on Windows/macOS &
  Delays CI and user installation &
  Provide pre-built wheels where possible; subprocess wrapper as fallback;
  Docker image for reproducible builds \\[4pt]

  R version incompatibility &
  \texttt{RobStatTM} may behave differently across R versions &
  Pin minimum R version (4.3+); test in CI against 4.3, 4.4, 4.5 \\[4pt]

  Numerical discrepancies in pure-Python ports &
  Undermines credibility of direct-porting approach &
  Tolerance tiers: $10^{-12}$ for wrappers, $10^{-6}$ for ports;
  investigate any discrepancy $> 10^{-4}$ before merging \\[4pt]

  Scope creep from PCRA extension &
  Core \texttt{robstatpy} quality suffers &
  PCRA is explicitly a stretch goal; only begin after
  Phase~4 benchmarks are complete and midterm evaluation is passed \\[4pt]

  Mentor availability during summer &
  Reduced feedback cadence &
  Three mentors provide redundancy; asynchronous communication
  via GitHub issues; weekly status updates regardless \\
  \bottomrule
\end{tabularx}
\renewcommand{\arraystretch}{1}

%--------------------------------------------------------------------
% SECTION 12 -- REFERENCES
%--------------------------------------------------------------------
\section{References}

\begin{enumerate}[leftmargin=2em, itemsep=3pt, label={[\arabic*]}]
  \item R.~A.~Maronna, R.~D.~Martin, V.~J.~Yohai, and
    M.~Salibian-Barrera.
    \textit{Robust Statistics: Theory and Methods (with R)}, 2nd ed.
    Wiley, 2019.
  \item M.~Salibian-Barrera, G.~Willems, and R.~Zamar.
    ``The fast-$\tau$ estimator for regression.''
    \textit{Journal of Computational and Graphical Statistics}, 2008.
  \item V.~J.~Yohai.
    ``High breakdown-point and high efficiency robust estimates for
    regression.''
    \textit{The Annals of Statistics}, 15(2):642--656, 1987.
  \item D.~L.~Rocke.
    ``Robustness properties of S-estimators of multivariate location and
    shape in high dimension.''
    \textit{The Annals of Statistics}, 24(3):1327--1345, 1996.
  \item L.~Gautier et al.
    \texttt{rpy2}: R in Python.
    \url{https://rpy2.github.io/}, 2024.
  \item R.~D.~Martin.
    \texttt{PCRA}: Portfolio Construction and Risk Analysis.
    CRAN, 2023.
\end{enumerate}

\end{document}
